\documentclass[10pt,a4paper]{amsart}
\usepackage[margin=19mm]{geometry}
\usepackage{amsmath,amssymb,mathtools}
\usepackage[scr=rsfs,cal=euler]{mathalfa}
\usepackage{xfrac}
\usepackage[unicode,hidelinks,bookmarks=false]{hyperref}
\newcommand{\ie}{i.e.\ }
\newcommand{\cf}{cf.\ }
\newcommand{\resp}{respectively\ }
\newcommand{\Subsection}{\subsection}
\providecommand{\nobreakdash}{\nobreak\hbox{-}}
\setlength{\emergencystretch}{2em}
\multlinegap0pt
\usepackage{amssymb}
\usepackage{mathtools}
\newtheorem{proposition}{Proposition}
\newtheorem{theorem}{Theorem}
\newcommand{\F}{\mathbf F}
\newcommand{\Fix}{\textnormal{Fix}}
\newcommand{\Q}{\mathbf Q}
\newcommand{\Z}{\mathbf Z}
\newcommand{\cycl}[1]{\Q(\zeta_{#1})}
\newcommand{\mc}{\mathcal}
\newcommand{\norm}[2]{N_{#1}^{#2}}
\title{Arithmetic exceptionality of generalized Chebyshev polynomials of the second kind}
\author{Derya Acar, Met\.{I}n Azmaz, Vural Cam, Ömer Küçüksakallı}
\date{Source: arXiv:2606.08125v1 (CC BY 4.0)}
\begin{document}
\maketitle

\noindent\emph{Source and licence.} Arithmetic exceptionality of generalized Chebyshev polynomials of the second kind. Version arXiv:2606.08125v1 (CC BY 4.0), distributed by arXiv under the Creative Commons Attribution 4.0 International licence, http://creativecommons.org/licenses/by/4.0/, as recorded in the arXiv licence metadata for this exact version and shown on the abstract page https://arxiv.org/abs/2606.08125v1. Full licence notice: \texttt{licenses/arxiv-2606.08125-CC-BY-4.0.txt}.\par\smallskip

\noindent\textit{Editorial scope.} Counted unit: the article's theorem thm:A1normQ (source0.tex lines 609--622). The article's complete original proof of it is delivered with it (source0.tex lines 623--657, one \texttt{proof} environment of 35 lines). No mathematics is written, repaired or re-derived: the statement and the proof are the article's own lines, in the article's own notation, and no retained line is substituted or re-flowed.  Retained uncounted as the source-local context the proof needs: lines 433--435; lines 437--443; lines 525--535; lines 547--549; lines 553--563.  Nothing was omitted from any retained span, so the package carries no omission record.  No retained line contains a citation, so the wrapper carries no bibliography and no bibliography entry.  No retained line contains a figure environment, an image-inclusion command or drawing code, so no image is delivered and the package declares no asset dependency.  Editorial formatting only, in addition: an AMS \texttt{amsart} preamble that restates the article's own declarations and macros used by the retained lines, namely the environments \texttt{proposition}, \texttt{theorem} on separate counters, together with the standard packages those lines need.  The retained environments print in this document as Proposition 1, Theorem 1 in order of appearance, whereas the article prints the same items with the numbering of its own full document.  The complete native source of the article as distributed in the arXiv e-print for this exact version is bundled unchanged as \texttt{sources/202307/source0.tex}.
\begin{equation}\label{eq:one-to-one}
	\F_q \leftrightarrow \Fix(P_{A_{1}}^{q}) = \mc{A}_{q} \cup \mc{B}_{q}.
\end{equation}

\begin{equation}\label{eq:QA1functionaltrigo}
	\begin{gathered}
		Q^k_{A_1}(2\cos{(2\pi u)})=\frac{\sin{(2\pi(k+1)u)}}{\sin{(2\pi u)}} \text{\quad for } u \in \left(0,\frac{1}{2}\right),\\
		Q^k_{A_1}(2)=k+1,\\
		Q^k_{A_1}(-2)=(-1)^k(k+1).
	\end{gathered}
\end{equation}

\begin{proposition}\label{prop:norm}
	Let $n$ and $j$ be positive integers such that $\gcd(j,n)=1$. Then,
	$$
	\norm{\Q}{\cycl{n}}\left(1-\zeta_n^j\right)=
	\begin{cases}
		0 & \text{if } n=1, \\
		l & \text{if } n=l^a \text{ for some } l\in\mathbb{P},  a \in \Z^{+}, \\
		1 & \text{otherwise}.
	\end{cases}
	$$
\end{proposition}

\begin{equation}\label{eq:norminextension}
N_{\Q}^{\cycl{d}}(\alpha_d) = \left( N_{\Q}^{\Q(\alpha_d)}(\alpha_d) \right)^2.	
\end{equation}

\begin{theorem}\label{thm:A1norminput}
The norm of $\alpha_d$ is given by:
\[N_{\Q}^{\Q(\alpha_d)}(\alpha_d) = 
\begin{cases} 
	2 & \text{if } d = 1, \\
	-2 & \text{if } d = 2, \\
	0 & \text{if } d = 4, \\
	\pm l & \text{if } d = 4l^a\text{ for some }l\in\mathbb{P}, a\in\Z^{+},\\
	\pm 1 & \text{otherwise}.
\end{cases}\]
\end{theorem}

\begin{theorem}\label{thm:A1normQ}
Let $k\ge 3$ be an odd integer and $q$ be a power of an odd prime $p$. Suppose that $p\nmid(k+1)$ and $$\gcd\left(\frac{q-1}{2}, k+1\right) \cdot \gcd\left(\frac{q+1}{2}, k+1\right) = 2.$$
If $d$ is a divisor of $q-1$ or $q+1$, then the norm of $Q_{A_1}^k(\alpha_d)$ is given by:
	$$
	N_{\Q}^{\Q(\alpha_d)}\left(Q_{A_1}^k(\alpha_d)\right) = 
	\begin{cases} 
		k+1 & \text{if } d = 1, \\
		(-1)^k(k+1) & \text{if } d = 2, \\
		0 & \text{if } d = 4, \\
		\pm l & \text{if } d = 4l^a\text{ for some }l\in\mathbb{P}, a\in\Z^{+},\\
		\pm 1 & \text{otherwise}.
	\end{cases}
	$$
\end{theorem}

\begin{proof}
The cases where $d$ equals $1$ or $2$ are obvious by (\ref{eq:QA1functionaltrigo}) as $\alpha_1=2$ and $\alpha_2=-2$, respectively. The case $d=4$ is also trivial as $\alpha_4=0$ is a zero of the polynomial $Q_{A_1}^k$ since $k$ is odd. For the remaining cases, we shall prove 
$$
N_{\Q}^{\Q(\zeta_d)}\left(Q_{A_1}^k(\alpha_d)\right) = 
\begin{cases}
	l^2 & \text{if } d = 4l^a\text{ for some }l\in\mathbb{P}, a\in\Z^{+},\\
	1 & \text{otherwise}.
\end{cases}
$$
in light of (\ref{eq:norminextension}). The functional equation that defines $Q_{A_1}^k$, namely \eqref{eq:QA1functionaltrigo}, allows us to write 
$$Q_{A_1}^k(\alpha_d)=\dfrac{\zeta_d^{k+1}-\zeta_d^{-(k+1)}}{\zeta_d-\zeta_d^{-1}}=\zeta_d^{-k}\left(\dfrac{1-\zeta_d^{2(k+1)}}{1-\zeta_d^2}\right).
$$
Using the multiplicativity of the norm and the fact that $\zeta_d^{-k}$ is a unit, we obtain the following equality
$$N_{\Q}^{\cycl{d}}(Q_{A_1}^k(\alpha_d))=\pm \frac{N_{\Q}^{\cycl{d}}\left(1-\zeta_d^{2(k+1)}\right)}{N_{\Q}^{\cycl{d}}\left(1-\zeta_d^2\right)}.
$$
We want to utilize Proposition~\ref{prop:norm}. We rely on the gcd-condition in the hypothesis to restrict our attention to rather limited cases. Depending on $q$ modulo $4$, the gcd-condition becomes one of the following
	$$
	\begin{cases}
		\gcd\left(\frac{q-1}{2}, k+1\right) = 2 \text{ and } \gcd\left(\frac{q+1}{2}, k+1\right) = 1 &\text{if } q\equiv1\pmod4,\\
		\gcd\left(\frac{q-1}{2}, k+1\right) = 1 \text{ and } \gcd\left(\frac{q+1}{2}, k+1\right) = 2 &\text{if } q\equiv3\pmod4.
	\end{cases}
	$$
As in the hypothesis, suppose that $d$ is a divisor of $q-1$ or $q+1$. We remark that this choice is compatible with the one-to-one correspondence \eqref{eq:one-to-one}, and we will eventually focus on $\alpha_d\in \Fix(P_{A_1}^q)$. 

We repeat the ideas in the proof of Theorem~\ref{thm:A1norminput}, emphasizing the differences between the expressions
\[ \frac{1-\zeta_d^4}{1-\zeta_d^2}, \quad \text{ and } \quad \dfrac{1-\zeta_d^{2(k+1)}}{1-\zeta_d^2}. \]

Suppose that $d=2^a$ with $a\ge 3$. Then the numerator comes with the power $2(k+1)$ instead of $4$. However, the gcd-condition implies that $1-\zeta_{2^a}^{2(k+1)}$ and  $1-\zeta_{2^{a-2}}$ are conjugates of each other. The rest of the proof is the same.  

Suppose that $d=4l^a$ for some odd prime $l$ with $a \ge 1$. Then the numerator is $1-\zeta_{4l^a}^{2(k+1)} = 1-\zeta_{l^a}^{(k+1)/2}$. Using Proposition~\ref{prop:norm} and the gcd-condition, we conclude that the numerator has norm $\pm l$ in the extension $\cycl{l^a}/\Q$. The rest of the proof is the same.

If $d>1$ is odd, then as $d$ is a divisor of $(q-1)/2$ or $(q+1)/2$, and the gcd-condition ensures that $d$ and $k+1$ are coprime. So, the numerator $1-\zeta_d^{2(k+1)}$ and the denominator $1-\zeta_d^2$ have the same norm as $1-\zeta_d$ by Proposition~\ref{prop:norm}.
	
If $d=2m$ with $m>1$ odd, then as $m$ is a divisor of $(q-1)/2$ or $(q+1)/2$, and the gcd-condition ensures that $m$ and $k+1$ are coprime. So, the numerator $1-\zeta_d^{2(k+1)}=1-\zeta_m^{k+1}$ and and the denominator $1-\zeta_d^4=1-\zeta_m^2$ have the same norm as of $1-\zeta_m$ by Proposition~\ref{prop:norm}.
\end{proof}

\end{document}
